On considère l’équation différentielle (E) : $y' + 3y = \left(2+2x\right)\e^{-2x}$ où $y$ est une fonction de la variable réelle $x$, définie et dérivable sur l’ensemble $\R$ des nombres réels, et $y'$ est la fonction dérivée de $y$.

\begin{enumerate}
	\item Donner et résoudre l'équation homogène (E$_0$)\sautp

	\Lignes{2}
    \item \begin{enumerate}
	\item A l’aide du logiciel Géogébra, déterminer les nombres réels $a$ et $b$ pour que la fonction $p$ définie sur $\R$ par : $p(x) = \left(a~x + b\right)\e^{-2x}$ soit une solution particulière de l'équation différentielle (E), et donner l’expression de $p$. \sautm
		
		%{\it On prendra des curseurs allant de $0$ à $40$ par pas de $2$}\sautm
		
		On trouve : $a=$\makebox[0.1\linewidth]{\dotfill} \hfill $b=$\makebox[0.1\linewidth]{\dotfill} \hfill $p(x)=$\makebox[0.3\linewidth]{\dotfill}\sauts
	\end{enumerate}
	\appelprof
    \begin{enumerate}
    \setcounter{enumii}{1}
	\item Vérifier que la fonction $p(x)=2x \e^{-2x}$ est une solution particulière de (E).\sautm

    \Lignes{5}
    \end{enumerate}
	\item En déduire la forme générale des solutions de (E).\sautp

	\Lignes{1}
    \item \begin{enumerate}
	\item Déterminer, à l'aide de géogebra et d'un curseur bien choisi, la solution particulière $f$ de (E) telle que $f'(0)=-1$.\sautp

	\Lignes{1}
    \end{enumerate}
    \appelprof
    \begin{enumerate}
    \setcounter{enumii}{1}
	\item Déterminer, par le calcul, la fonction $f$ solution de (E) telle que $f'(0)=-1$.\sautp

    \Lignes{5}
    \end{enumerate}
\end{enumerate}